\chapter*{Hướng Dẫn Sử Dụng Sách}
\addcontentsline{toc}{chapter}{Hướng Dẫn Sử Dụng Sách}
\markboth{Hướng Dẫn Sử Dụng Sách}{Hướng Dẫn Sử Dụng Sách}

\begin{truyen}{Cách Đọc Cuốn Này}{How to Read This Book}
\chuhoa{C}{uốn này dày hơn nên không nên đọc liền một mạch như đọc tin ngắn. Tốt hơn là đọc theo từng phần: vài ngày cho \textbf{Được và Mất}, vài ngày cho \textbf{Thắng và Thua}, rồi chậm lại ở \textbf{Đúng và Sai}.}
\textit{(This book is thicker and should not be skimmed like short news. It is better read by section: a few days for \textbf{Gain and Loss}, a few for \textbf{Winning and Losing}, then a slower pace for \textbf{Right and Wrong}.)}

\textbf{Người dạy học:} có thể dùng mỗi phần như một cụm chủ đề riêng để kể chuyện hoặc gợi thảo luận.
\textit{(\textbf{Teachers:} may use each section as its own theme cluster for storytelling or discussion.)}

\textbf{Người đọc cá nhân:} nên đọc 1--2 truyện một ngày, đánh dấu câu nào chạm mình nhất.
\textit{(\textbf{Individual readers:} should read 1--2 stories a day and mark the lines that touch them most.)}
\end{truyen}

\begin{baihoc}
Cuốn sách này hiệu quả nhất khi đọc chậm theo từng phần, chứ không phải đọc nhanh cho xong số trang.
\textit{(This book works best when read slowly by section, not quickly just to finish the page count.)}
\end{baihoc}

\begin{ghinhoanh}
\textbf{Reading tips:}

Read by section, not by hurry.

One honest reflection is worth more than many rushed pages.
\end{ghinhoanh}

\ngancach